\chapter{DISCUSSION ET PERSPECTIVES}

\section{Analyse critique du projet}

\subsection*{Ce qui a bien fonctionné}

Le résultat le plus concret est simple : les trois pipelines tournent, les sorties PTF et Capitaux sont identiques au SAS, et la documentation existait. Aucun de ces trois points n'était acquis au départ.

Le choix d'\textbf{analyser avant de coder} a été le bon. Deux semaines investies dans l'analyse du code SAS ont évité des erreurs d'interprétation qui auraient pu passer inaperçues jusqu'à la validation finale. C'est contre-intuitif pour un développeur, mais sur un projet de migration sans documentation initiale, c'est indispensable.

L'\textbf{architecture médaillon} a aussi montré son utilité pendant la validation : chaque couche étant persistée, il a été possible d'inspecter les données à chaque étape et d'isoler précisément l'origine d'une divergence. Sans cette séparation, déboguer la correction \texttt{cdpole} aurait été bien plus laborieux.

La \textbf{configuration externalisée} s'est révélée très pratique à l'usage. Corriger un filtre d'exclusion ou ajouter un code produit n'implique pas de toucher au code Python, juste modifier un fichier JSON. C'est une vraie différence par rapport au SAS où les paramètres étaient éparpillés dans le code.

\subsection*{Ce qui a été difficile}

La principale contrainte a été \textbf{l'absence totale de documentation SAS}. Toute la logique de traitement résidait dans le code, dans des macros imbriquées et des noms de variables peu expressifs. Cela a ralenti le démarrage et a créé un risque résiduel : malgré le soin apporté à l'analyse, il est impossible d'exclure complètement des règles implicites qui auraient pu passer inaperçues.

L'\textbf{accès aux données de test} a aussi été compliqué. Les protocoles de sécurité Azure d'Allianz sont stricts, et plusieurs transformations ont dû être développées sans données réelles. Le délai entre une demande d'accès et sa réalisation effective a parfois allongé les cycles de correction de plusieurs jours. C'est une contrainte inhérente au contexte industriel, mais elle doit être anticipée dans le planning des migrations futures.

\section{Positionnement par rapport aux pratiques de l'industrie}

La solution s'aligne bien avec les standards actuels du data engineering. L'architecture médaillon est aujourd'hui la recommandation centrale de Databricks pour les projets lakehouse, et le choix d'Azure Databricks correspond à la stratégie cloud du groupe Allianz, qui a standardisé ses infrastructures sur Azure depuis 2022.

Par rapport à d'autres approches observées dans l'assurance, notamment les migrations vers Pandas qui atteignent leurs limites sur des volumes de quelques gigaoctets, PySpark garantit la capacité à traiter des volumes bien supérieurs au datamart Construction actuel sans refonte. Le respect des conventions PEP 8, des types Python et d'une structure modulaire représente aussi une rupture nette avec le code SAS d'origine, qui ne suivait aucune convention formelle.

\section{Ce que ça apporte à Allianz France}

Trois gains concrets sont attendus :

\begin{itemize}
    \item \textbf{Réduction des coûts de licences} : SAS est un logiciel propriétaire coûteux. Le passage à une stack Python open source sur une infrastructure Azure déjà en place supprime ce poste de dépense récurrent.
    \item \textbf{Scalabilité} : là où le serveur SAS est une ressource à capacité fixe, un cluster Spark sur Azure s'adaptte dynamiquement au volume à traiter. L'ajout de nouvelles visions historiques ou l'augmentation du périmètre n'impliquent aucun investissement supplémentaire en infrastructure.
    \item \textbf{Capitalisation documentaire} : avant ce projet, la connaissance du datamart Construction était uniquement contenue dans un code non maintenu. Elle est maintenant dans plusieurs documents accessibles à toute l'équipe. Cet actif est durable, indépendamment du sort du projet de migration.
\end{itemize}

Deux points de vigilance méritent d'être mentionnés. Le \textbf{double environnement}, SAS et Python coexistant tant que la validation n'est pas complète, a un coût opérationnel qui doit être provisionné. Et la \textbf{montée en compétences PySpark} pour les profils historiquement SAS demande un accompagnement qui ne peut pas être improvisé.

\section{Perspectives}

\textbf{Court terme} : finaliser la validation de parité sur dix visions historiques pour les pipelines Capitaux et Émissions. Une fois cette étape franchie, le pipeline PTF Mouvements est prêt pour une mise en production.

\textbf{Moyen terme} : le framework développé, module de lecture, logger, configuration JSON, structure de projet, est directement réutilisable pour les quatre marchés suivants (Auto, Habitation, Santé, Épargne). Les prochaines migrations pourront se concentrer sur les règles métier spécifiques à chaque marché, sans repartir de zéro pour les briques communes. Quatre recommandations découlent de l'expérience de ce premier prototype :

\begin{enumerate}
    \item Commencer par la documentation fonctionnelle avant tout développement.
    \item Intégrer les délais d'accès aux données Azure dans le planning prévisionnel.
    \item Réutiliser les modules génériques sans les réécrire.
    \item Organiser des revues de parité intermédiaires tout au long du développement.
\end{enumerate}

\textbf{Long terme} : quand tous les marchés seront migrés et validés, l'infrastructure SAS pourra être désengagée. L'évolution vers Delta Lake, format ACID avec \textit{time travel}, est une amélioration naturelle à envisager pour les marchés suivants.
